\section{Resultados II --- Síntese por janela}\label{cap:janelas}

Este capítulo lê o corpus longitudinalmente, nas três janelas tecnológicas. A
Tabela~\ref{tab:sintese_janelas} resume, por janela, a tecnologia e o tipo de
estudo dominantes, o sinal e a polarização predominantes e a fração de estudos
que mobiliza cada mecanismo teórico do arcabouço de tarefas de Acemoglu e
Restrepo --- deslocamento, reinstalação, complementaridade e demanda agregada.%
\footnote{Os mecanismos não são mutuamente exclusivos: um mesmo estudo pode
invocar mais de um, de modo que os percentuais somam mais de 100\%.}
A Figura~\ref{fig:mecanismos_janela} traça a evolução desses mecanismos.

Dois traços atravessam as três janelas e antecipam a discussão de continuidades
do Capítulo~\ref{cap:comparacao}: em todas elas o \emph{sinal} predominante
sobre o emprego é ``ambíguo'', e nas duas janelas de tamanho analisável a
\emph{polarização} modal é o risco concentrado em ocupações de baixa
qualificação. A literatura econômica, vista no agregado, nunca convergiu para um
veredito unívoco sobre o saldo líquido de empregos, e a tese da polarização ---
risco maior para tarefas rotineiras de menor qualificação --- mantém-se como
leitura mais frequente ao longo de todo o período. O que muda entre as janelas é
o enquadramento tecnológico, o tipo de evidência e o peso relativo dos
mecanismos.

\begin{figure}[htbp]
\centering
\includegraphics[width=0.9\linewidth]{figures/mecanismos_janela.pdf}
\caption{Mecanismos teóricos invocados por janela temporal.}
\label{fig:mecanismos_janela}
\end{figure}

\begin{table}[htbp]
\centering
\caption{Síntese das dimensões por janela temporal.}
\label{tab:sintese_janelas}
{\footnotesize\input{tables/sintese_janelas.tex}}
\end{table}

\subsection{Janela 1: \janelaUm{} --- era da automação}

A primeira janela é muito pequena (15 estudos) e serve apenas de pano de fundo:
qualquer percentual calculado sobre ela muda vários pontos com a reclassificação
de um único estudo. Ainda assim é nítida no enquadramento. A tecnologia dominante
é a automação, e o mecanismo de \emph{deslocamento} é onipresente --- todos os
estudos que classificaram a dimensão o invocam ---, enquanto a \emph{demanda
agregada} comparece de forma marginal (13,3\%). O período retrata a tecnologia
pela lente do efeito-substituição clássico: a máquina desloca tarefas rotineiras,
e o debate gira em torno de quanto desse deslocamento é compensado por
reinstalação (46,7\%) e complementaridade (66,7\%). São representativos da
janela a decomposição entre comércio e tecnologia nos mercados de trabalho locais
de \textcite{autorDornHanson2015} e a revisão crítica das estimativas de risco de
automação de \textcite{arntzGregoryZierahn2017}.

\subsection{Janela 2: \janelaDois{} --- era de deep learning e ML}

Na segunda janela (225 estudos) a automação segue como tecnologia dominante e o
tipo modal de estudo é o teórico ou de modelagem, seguido de perto pela evidência
macro/setorial: é o período em que o arcabouço de tarefas é formalizado e
estendido. O deslocamento permanece quase universal (95,1\%) e a \emph{demanda
agregada} atinge seu ápice no corpus (24,2\%): a difusão dos modelos preditivos
convive com um debate explícito, em boa parte de equilíbrio geral, sobre os
efeitos de produtividade e de expansão da demanda que poderiam recompor o
emprego, como na análise do ajuste dos mercados de trabalho à robótica industrial
de \textcite{dauth2021}. A reinstalação (47,9\%) e a complementaridade (57,1\%)
são invocadas por cerca de metade dos estudos. É nesta janela que a polarização
com risco na \emph{baixa} qualificação é mais hegemônica: apenas 4 dos 138
estudos que classificaram a dimensão situam o risco na alta qualificação.

\subsection{Janela 3: \janelaTres{} --- era da IA generativa}

A terceira janela (576 estudos) marca três inflexões. Primeiro, no tipo de
evidência: a modelagem teórica perde espaço (de 29,3\% para 13,2\% dos estudos)
e o estudo modal passa a ser a evidência macro/setorial, ao lado de um forte
crescimento das revisões e dos estudos de firma. Segundo, nos mecanismos: a
\emph{complementaridade} salta para 73,0\% --- o maior valor entre as janelas de
tamanho analisável ---, enquanto o deslocamento recua dez pontos (85,1\%) e a
demanda agregada cai pela metade (12,4\%). A literatura da era generativa, sem
abandonar a preocupação com a substituição, enfatiza de forma crescente a
ampliação de capacidades e a cooperação humano-máquina, a exemplo da evidência
de firma sobre IA generativa de \textcite{brynjolfssonLiRaymond2025}; ao mesmo
tempo, torna-se mais microeconômica e menos atenta aos canais de equilíbrio
geral. Terceiro, a rotulagem tecnológica torna-se mais difusa (``geral'' passa a
dominar), sinal de que muitos trabalhos tratam a IA de modo amplo mesmo quando o
gatilho é a onda generativa. Essas inflexões são submetidas a teste formal na
comparação pareada do próximo capítulo. Convém antecipar uma cautela: a terceira
janela é definida pelo calendário, e não pelo conteúdo, e menos de um décimo de
seus estudos tem a IA generativa como objeto. O que a janela registra é a
literatura \emph{contemporânea} da IA generativa, em boa parte ainda dedicada à
automação e à robótica; a Discussão (Capítulo~\ref{cap:discussao}) separa as duas
coisas.
